\chapter{课程地图、工具与学习方式}

\section{一条消息怎样经过系统}
先从一笔订单开始：用户付款后，商店要把“订单已付款”交给后续处理程序，例如库存服务。若付款程序直接调用库存程序，接收方停机时消息可能丢失；消息系统让发送方先交出一条消息，由系统保存，再由接收方按自己的速度读取。Apache Kafka 是一种分布式消息流平台；本书用可运行的小型教学系统解释它所涉及的基本问题，不假定读者已知道 Kafka，也不把课程协议说成官方 Kafka 协议。

本书把一条可保存、可读取的数据称为\textbf{消息}或\textbf{记录（record）}。\textbf{生产者（producer）}发送记录，\textbf{broker}接收请求并保存记录，\textbf{消费者（consumer）}主动请求并处理记录。一个消息系统是负责在它们之间传递和保存记录的软件。发送请求成功、记录被消费者读到、业务处理完成，是三个不同的事件；后面的章节会分别保存和检查它们的状态。

\textbf{topic（主题）}是给相关消息取的分类名称，例如 \pathcode{orders}。一个 topic 可以拆成多条各自有序的\textbf{partition（分区）}日志；本课程用分区让不同记录可以独立追加和读取。记录的 \textbf{key} 可用于关联或路由，\textbf{value} 是实际内容，\textbf{timestamp（时间戳）}是随记录保存的时间值；时间戳不替代决定分区内顺序的 \textbf{offset（位点）}。每条记录的 offset 是它在所属分区里的编号，从 0 开始；不同分区各自编号，不能跨分区比较，也不存在一个代表整个 topic 的全局顺序。
\begin{center}
\small
\begin{tabular}{lll}
\toprule
topic & partition & 记录的 offset 顺序 \\
\midrule
\pathcode{orders} & 0 & 0，1 \\
\pathcode{orders} & 1 & 0，1 \\
\bottomrule
\end{tabular}
\end{center}
表中两个分区的 offset 0 是两条不同记录；本书能保证各分区内部的顺序，不给两个分区之间排出先后。

一条记录的生命周期是：生产者选择目标分区；broker 把记录追加到该分区日志；消费者主动拉取记录；程序执行处理回调；处理成功后再保存消费进度。读取不会删除日志中的记录，所以另一消费者或重启后的消费者仍可能再次读取它。日志保留决定旧记录何时被删；消费进度只说明某个消费者从哪里继续，两者不是同一件事。第 4 章会解释处理成功与保存进度之间的失败窗口。

\begin{center}
\small
\begin{tabular}{p{.22\textwidth}p{.68\textwidth}}
\toprule
章节与步骤 & 要解决的问题 \\
\midrule
第 1 章，1--4 & 怎样编码记录、追加和读取文件，并识别崩溃留下的不完整尾部？\\
第 2 章，5--8 & 怎样用索引加速查找、分段保存，并分别删除旧头部或截断尾部？\\
第 3 章，9--12 & 怎样选择分区、发现 broker 的位置，并把 TCP 字节还原成完整请求？\\
第 4 章，13--16 & 怎样批量发送、读取记录、保存消费进度，以及理解失败重放？\\
第 5 章，17--20 & 多个消费者如何分配分区，旧成员的请求如何被拒绝？\\
第 6 章，21--24 & broker 间怎样复制日志，哪些记录可以被普通消费者看见？\\
第 7 章，25--28 & leader 故障时怎样选择副本、修复分叉并重新接纳副本？\\
第 8 章，29--31 & 如何观察跨层故障、识别重复发送，并让本地业务效果可重放？\\
\bottomrule
\end{tabular}
\end{center}

\section{本课程实现什么（以及不是什么）}
本项目把上述路径做成一个小而可运行的系统：磁盘上的分区日志、TCP 请求、生产与消费、消费位点、单 topic 消费组，以及多个本机 broker 之间的拉取复制和受控切换。它不是 Kafka 官方客户端，也不使用 Kafka wire protocol；不能把官方 \method{KafkaProducer} 或 \method{KafkaConsumer} 指向本程序。

三个 Java source root 是课程设计的一部分。\pathcode{provided/} 放数据类型、I/O、传输、序列化样板、生命周期和测试夹具；它让注意力集中在算法而非重复搭框架。\pathcode{exercises/} 是学生唯一需要逐步补全的工作区，所有步骤的类型和公开签名预先可见，尚未到达的方法会抛出带步骤号的 \method{ExerciseNotImplementedException}。\pathcode{reference/} 是隔离的完整实现，供尝试后对照契约和课程验证；它与练习代码类名相同，因此学生和参考源根绝不能同次编译。测试在两种模式下复用同一套黑盒契约；学生运行时 classpath 不含参考 classes。

这不是生产级 Kafka：没有 Kafka 线协议兼容、事务、压缩、compaction、完整控制器共识或 controller HA。第 8 章的有限幂等生产者协议和本地持久化效果账本不是 Kafka 生产协议，也不构成端到端事务。它们使用单 JVM 内可信的教学控制面；虽有多台 broker 的 TCP server、线程和独立磁盘目录，副本仍共享该控制面。它能演示受控故障切换，不能证明任意网络分区下具备共识或生产级可用性。具体差异会在相关步骤说明。

\section{准备 Java 21 与课程命令}
项目使用 Gradle Wrapper \coursecmd{./gradlew}；Gradle 9.8.1 可运行于 Java 17+，Java 编译、测试和 demo 任务需要 Java 21+ JDK，并以 \pathcode{--release 21} 编译。不需要 Maven、Docker 或已启动的 Kafka。脚本使用 Bash 3.2+、\pathcode{find}、\pathcode{awk}、\pathcode{curl} 与 \pathcode{sha256sum} 或 \pathcode{shasum}；准备 Tectonic 需要 \pathcode{tar} 和 \pathcode{install}，macOS 解压 Pandoc 还需要 \pathcode{unzip}。EPUB 图形需要 Poppler 的 \pathcode{pdftoppm}；\coursecmd{./gradlew :setupBook} 会检查它，但不安装系统软件。Linux/WSL 使用 \coursecmd{sudo apt-get install poppler-utils}，macOS 使用 \coursecmd{brew install poppler} 安装 Poppler。先检查 \coursecmd{java -version}、\coursecmd{javac -version} 与 \coursecmd{./gradlew --version}。可通过 \pathcode{JAVA_HOME} 或 \pathcode{PATH} 使用任意符合要求的 JDK；SDKMAN 可选，项目不会静默修改 SDKMAN、系统软件或 shell 配置。

\coursecmd{./gradlew :setup} 通过 Gradle 解析固定版本 JUnit 1.11.4 并验证 checksum，不再把 JUnit 放入 \pathcode{.tools/}。\coursecmd{./gradlew :setupBook} 准备经过 checksum 校验的 Tectonic 0.17.0 与 Pandoc 3.12.1 到 \pathcode{.tools/}，并检查 Poppler；不安装系统软件。\coursecmd{./gradlew :doctor} 检查 Gradle、Java、Gradle 管理的 JUnit 与可选教材工具。\coursecmd{./gradlew :compile} 编译学生骨架与测试；编译成功不表示练习已经实现。课程清单与教材默认英文；设置 \pathcode{COURSE_LANG=zh} 后运行 \coursecmd{./gradlew :list} 可显示中文步骤标题。\coursecmd{./gradlew :book} 编译英文 PDF/EPUB 到 \pathcode{build/book/en/simpleKafka.pdf} 与 \pathcode{build/book/en/simpleKafka.epub}；设置 \pathcode{COURSE_LANG=zh} 后运行同一命令，生成 \pathcode{build/book/zh/simpleKafka-zh.pdf} 与 \pathcode{build/book/zh/simpleKafka-zh.epub}。每次还会更新所选语言的 \pathcode{docs/book/dist/} 分发文件；读者可直接打开已分发成品，无需 Java 或编译工具。book 任务需要启动 Wrapper JVM，但不编译 Java 或解析 JUnit。首次构建需联网获取 TeX bundle；两种语言都在线成功后，设置 \pathcode{BOOK_OFFLINE=1} 并用 \coursecmd{./gradlew :book --offline} 离线构建。Gradle \pathcode{--offline} 控制 Gradle 依赖；\pathcode{BOOK_OFFLINE=1} 为两次 Tectonic 调用启用 TeX bundle only-cached 模式。

\section{读懂二进制与文件 I/O}
文件只保存字节，不会替程序记住字段边界。UTF-8 字符“值”占 3 个字节，因此课程例子中的 value 长度是 3，而不是字符个数 1；普通记录的 length=28+3=31，总占用=4+31=35 字节。编码格式必须同时规定字段顺序、整数宽度和字节序。本课程用大端序：例如整数 1 的四个字节是 \pathcode{00 00 00 01}，最高有效字节在前。

\method{ByteBuffer} 是一段带容量、position 与 limit 的字节视图。position 指下一次读写的位置；limit 是当前可读/可写范围的末端。写入模式下调用 \method{flip()} 会把 position 归零、limit 设为刚写完的位置，随后便可读取这些字节；不能把 flip 当成清空数据。buffer 的 position 是其中的字节游标，后续练习中的 record offset 是分区内的逻辑编号，不能互换。解码器只应推进到已成功消费的边界。

\method{FileChannel} 可以从指定文件 byte position 读取或写入，也能请求强制同步已写数据。一次 \method{channel.read} 可能只填入部分 buffer，一次 \method{write} 也可能只写出部分内容；课程提供的 \method{ChannelIO.readFully/writeFully} 会循环处理这种情况。步骤 2 会用锁（lock）让同一日志上的并发追加一个接一个完成，避免批次交错；锁只保护短暂的共享状态操作，不应拿着它等待网络。每批追加后的 \method{force(false)} 是请求持久化文件数据的简化方式，不是对掉电、电源缓存、介质损坏或文件系统语义的数学保证。CRC32C 检测记录字节意外改变，但不修复错误，也不能证明数据来源可信。

\section{测试失败的阅读法}
先读失败的 Step 测试名和断言，再回到本章契约。测试报告中的第一处异常通常比最后一段堆栈更有用：未实现方法给出步骤和方法名；课程错误应保留 \method{ErrorCode}；意外 IOException 可能被包装为 \method{STORAGE_ERROR} 并保留 cause。不要把错误映射成空列表或成功默认值来让测试变绿。一个单步定位命令只运行该步，不代表累计依赖都满足；累计命令会再次评估之前每步。
课程 runner 按步骤分组输出测试名称、各步和总计数；失败时显示断言、精确方法定位与重跑命令。报告始终为英文；legacy JUnit XML 保存到 \pathcode{build/student/reports/} 或 \pathcode{build/reference/reports/}。自定义 JavaExec 课程任务输出纯文本。标准 Gradle \pathcode{:exercises:test} 与 \pathcode{:reference:test} 支持 IDE 调试及精确 \pathcode{--tests} 选择；其 XML/HTML 报告位于 \pathcode{build/<mode>/gradle-test-results/} 与 \pathcode{build/<mode>/gradle-test-reports/}。runner 保留 \pathcode{NO_COLOR} 与 \pathcode{TERM=dumb} 颜色门禁。

\begin{center}
\small
\begin{tabular}{p{.48\textwidth}p{.43\textwidth}}
\toprule
命令 & 意义 \\
\midrule
\coursecmd{./gradlew :list} & 显示 31 步、编辑方法与前置关系。\\
\coursecmd{./gradlew :stepTest -Pstep=12} & 只定位第 12 步；失败时辅助找当前问题。\\
\coursecmd{./gradlew :test -Pstep=12} & 累计运行 Step01Test 至 Step12Test。\\
\coursecmd{./gradlew :referenceTest -Pstep=31} & 对隔离参考实现运行完整 31 步契约。\\
\coursecmd{./gradlew :referenceDemo} & 运行完整参考 TCP/failover 场景。\\
\coursecmd{./gradlew :referenceConsistencyDemo} & 运行第 8 章真实 TCP 与磁盘一致性场景。\\
\bottomrule
\end{tabular}
\end{center}

步骤编号接受 1–31 的十进制数字，可有前导零（例如 \pathcode{-Pstep=1} 或 \pathcode{-Pstep=01}）。章节边界是第 4、8、12、16、20、24、28、31 步：只有该边界累计测试全绿，才算章节完成。\coursecmd{./gradlew :test -Pstep=8} 通过意味着第 1–8 步全部通过；\coursecmd{./gradlew :stepTest -Pstep=8} 仅定位第 8 步，不授予晋级。全课程测试不应把参考源放进学生 classpath，参考模式也不能代替自己的练习进度。

\section{已学概念小结与后续导航}
先记住三点：主题是记录分类，分区是各自有序的日志，而每条记录的 offset 只在自己的分区内有意义；发送、读取和业务处理是不同事件，读取本身不删除记录；日志保留决定记录何时删除，消费进度决定某个消费者从哪里继续。后续每遇到新的术语和不变量，都会先在对应步骤解释，不需要预先掌握 Kafka 缩写。

可以带着具体问题进入后续步骤：\stepref{1}{记录编解码}解释一条记录如何变成有明确边界的字节；\stepref{4}{崩溃恢复}说明如何判断文件末尾是未写完还是完整数据损坏；\stepref{14}{拉取消费者}展示消费者读取位置如何推进；\stepref{22}{ISR与高水位}解释多份日志进度不一致时，普通消费者能看到哪些记录；\stepref{25}{干净选举与epoch}讲一个 broker 接替写入角色后，如何拒绝旧角色发来的请求。
